
\label{sec:cilindricas}
\vspace{0.2cm}

Para el cálculo de integrales en regiones con simetría cilíndrica, se utiliza el sistema de coordenadas cilíndricas. Un punto $(x, y, z)$ se relaciona con $(r, \theta, z)$ mediante:

\begin{equation}\label{def:coord_cilindricas}
\mathbf{T}:W^*=[0,R]\times[0, 2\pi]\times[0,H] \subseteq\Rn{3} \to W\: : \: \mathbf{T}(r, \theta,z ) =( r \cos \theta , r \sin \theta,z ).
\end{equation} 

\begin{tarea} Probar que la transformaci\'on $T$ definida en (\ref{def:coord_cilindricas}) est\'a bajo las las hi\'otesis del Teorema  ~\ref{teo:cambio_variable}
\end{tarea}

\vspace{0.2cm}

Para el c\'alculo del \emph{Determinante Jacobiano} de $T$  primero  calculamos su  matriz diferencial $D_\mathbf{T}$, quedando
       \begin{equation*}
       \boldsymbol{D}_{\mathbf{T}}(r, \theta,z ) 
       =
       \begin{pmatrix}
       \cos\theta & -r\sin\theta & 0 \\
       \sin\theta & r\cos\theta & 0 \\
       0 & 0 & 1
       \end{pmatrix}
 \end{equation*}
 luego, el jacobiano es
       \begin{equation}\label{ejemplo_jacobiano_polares}
       J_{\mathbf{T}}(r, \theta,z) = \det(\boldsymbol{D}_{\mathbf{T}})(r, \theta,z) = 
      r(\cos^2\theta + \sin^2\theta) = r,
      \end{equation} 
 por \'ultimo,  en nuestro caso,  como $r \geq 0$, tenemos
\begin{equation}
|J_{\mathbf{T}}| = r
\end{equation}

\textbf{Visualización del sistema de coordenadas cilíndricas}

\begin{figure}[H]
    \centering
    \tikzsetnextfilename{cilindricas_cambio_figura}

    \begin{tikzpicture}

    %%% IZQUIERDA: caja W* en espacio (r, θ, z)
    \begin{axis}[
        name=izq,
        view={110}{25},
        width=7cm,
        height=7cm,
        xmin=-0.2, xmax=3,
        ymin=-0.2, ymax=7.5,
        zmin=-0.2, zmax=3,
        xlabel={$r$},
        ylabel={$\theta$},
        zlabel={$z$},
        axis lines=center,
        ticks=none,
        every axis x label/.style={at={(ticklabel* cs:1.05)}, anchor=north east},
        every axis y label/.style={at={(ticklabel* cs:1.05)}, anchor=north west},
        every axis z label/.style={at={(ticklabel* cs:1.05)}, anchor=south},
    ]

    % Cara inferior (z = 0)
    \addplot3[fill=magenta!20, opacity=0.85, draw=magenta!70, thick]
        coordinates {(0,0,0)(2.5,0,0)(2.5,6.28,0)(0,6.28,0)(0,0,0)};
    % Cara posterior (θ = 0)
    \addplot3[fill=magenta!15, opacity=0.7, draw=magenta!70, thick]
        coordinates {(0,0,0)(2.5,0,0)(2.5,0,2.5)(0,0,2.5)(0,0,0)};
    % Cara lateral (r = 0)
    \addplot3[fill=magenta!20, opacity=0.7, draw=magenta!70, thick]
        coordinates {(0,0,0)(0,6.28,0)(0,6.28,2.5)(0,0,2.5)(0,0,0)};
    % Cara frontal (r = 2.5)
    \addplot3[fill=magenta!30, opacity=0.8, draw=magenta!70, thick]
        coordinates {(2.5,0,0)(2.5,6.28,0)(2.5,6.28,2.5)(2.5,0,2.5)(2.5,0,0)};
    % Cara lateral (θ = 2π)
    \addplot3[fill=magenta!25, opacity=0.75, draw=magenta!70, thick]
        coordinates {(0,6.28,0)(2.5,6.28,0)(2.5,6.28,2.5)(0,6.28,2.5)(0,6.28,0)};
    % Cara superior (z = 2.5)
    \addplot3[fill=magenta!35, opacity=0.9, draw=magenta!70, thick]
        coordinates {(0,0,2.5)(2.5,0,2.5)(2.5,6.28,2.5)(0,6.28,2.5)(0,0,2.5)};

    % Etiqueta W*
    \node[color=magenta!80!black, font=\Large]
        at (axis cs: 1.25, 3.14, 1.25) {$W^*$};

    % Cotas
    \node[below, color=black!60, font=\small]
       at (axis cs: 2.5, 0, 0) {$R$};
    \node[above right, color=black!60, font=\small]
        at (axis cs: 0, 6.28, 0) {$2\pi$};
    \node[left, color=black!60, font=\small]
        at (axis cs: 0, 0, 0) {$z_1$};
    \node[left, color=black!60, font=\small]
        at (axis cs: 0, 0, 2.5) {$H$};

    \end{axis}

    % Label inferior izquierdo
    \node[below=0.15cm of izq.south, font=\small]
        {Espacio $(r,\theta,z)$};

    %%% FLECHA T
    \draw[-{Stealth[length=9pt, width=6pt]}, thick]
        ($(izq.east)+(0.4,0)$)
        -- node[above, font=\large]{$T$}
        ++(1.5,0);

    %%% DERECHA: cilindro W en espacio (x, y, z)
    \begin{axis}[
        name=der,
        at={($(izq.east)+(2.4cm,0)$)},
        anchor=west,
        view={110}{25},
        width=7cm,
        height=7cm,
        xmin=-2.8, xmax=2.8,
        ymin=-2.8, ymax=2.8,
        zmin=-0.2, zmax=3,
        xlabel={$x$},
        ylabel={$y$},
        zlabel={$z$},
        axis lines=center,
        ticks=none,
        every axis x label/.style={at={(ticklabel* cs:1.05)}, anchor=north east},
        every axis y label/.style={at={(ticklabel* cs:1.05)}, anchor=north west},
        every axis z label/.style={at={(ticklabel* cs:1.05)}, anchor=south},
    ]

    % Disco inferior (z = 0)
    \addplot3[surf, fill=green!25, draw=green!55,
              fill opacity=0.8, draw opacity=0.9,
              domain=0:360, y domain=0:2,
              samples=36, samples y=3, shader=flat]
        ({y*cos(x)}, {y*sin(x)}, {0});

    % Superficie lateral
    \addplot3[surf, fill=green!20, draw=green!45,
              fill opacity=0.65, draw opacity=0.8,
              domain=0:360, y domain=0:2.5,
              samples=36, samples y=10, shader=flat]
        ({2*cos(x)}, {2*sin(x)}, {y});

    % Disco superior (z = 2.5)
    \addplot3[surf, fill=green!35, draw=green!60,
              fill opacity=0.9, draw opacity=1,
              domain=0:360, y domain=0:2,
              samples=36, samples y=3, shader=flat]
        ({y*cos(x)}, {y*sin(x)}, {2.5});

    % Etiqueta W
    \node[color=green!50!black, font=\Large]
        at (axis cs: 0.5, 0.5, 1.25) {$W$};

    \end{axis}

    % Label inferior derecho
    \node[below=0.15cm of der.south, font=\small]
        {Espacio $(x,y,z)$};

    \end{tikzpicture}

    \caption{Visualización del sistema de coordenadas cilíndricas. 
    El bloque $W^*$ a la izquierda representa la región en el espacio 
    $(r,\theta,z)$ con sus límites típicos, mientras que el cilindro 
    $W$ a la derecha muestra cómo se transforma esa región en el 
    espacio $(x,y,z)$ mediante la transformación $T$. \href{https://www.geogebra.org/m/yam9nnxq}{Ver en GeoGebra }}

\end{figure}